% -----------------------------------------------------------------------------
%  02_method_baseline.tex
%  Every numeric value here is a macro from macros.tex (authority: NUMBERS.md).
%  Values with no macro are marked \needsnum{...} and must be resolved before
%  submission -- do not replace them with literals read out of results/*.json.
% -----------------------------------------------------------------------------
\section{Method: how not to build a strawman}
\label{sec:method}

The comparison in this paper is between architectures, but what is actually
measured is a pair of implementations. Everything downstream --- the cost laws,
the walls, the mechanism --- rests on the claim that the two sides were priced
equally. This section is the audit of that claim. It is also, we think, the most
transferable part of the work: both defects we found were invisible in the
result tables, both were defaults rather than decisions, and both flattered the
architecture under test.

\subsection{Experimental setting}

All numbers in this section come from a single campaign: $n = 2000$ BGL
(LOGHUB) events from the temporal test split, a fixed 12-field extraction
schema, BF16, an H100 PCIe 80GB host, and vLLM 0.19.1 for the served
configurations. The HuggingFace-transformers configurations that these replace
ran the \emph{same} events with the \emph{same} prompts, so quality is directly
comparable across serving stacks and any difference between them is a property
of the stack, not of the data. Greedy decoding throughout; no grammar
constraint, no constrained decoding, on either side.

\subsection{Defect 1: a decoding default that penalises copying}
\label{sec:defect1}

Qwen2.5 ships a \texttt{repetition\_penalty} in its \texttt{generation\_config.json}:
\repPenaltySmall\ at 1.5B and \repPenaltyLarge\ at 3B and 7B. HuggingFace
\texttt{generate()} honours that penalty \emph{even under greedy decoding}
(\texttt{do\_sample=False}), because the penalty is a logits processor and the
sampling flag only selects the reduction applied afterwards. Our AR wrapper
never set the parameter, so the checkpoint's own default silently took effect on
every AR run in the study.

For this task the default is not merely suboptimal, it is inverted. Correct
extraction consists largely of copying spans verbatim out of the input --- host
names, component identifiers, severities, timestamps --- and a repetition
penalty down-weights exactly the tokens already present in the context. The
signature is unmistakable in the per-field breakdown: only copy-from-input
fields moved when we neutralised the penalty, while the fields the model has to
infer were untouched. The most affected field was \texttt{host}, emptied on
\hostEmptied\ events. Neutralising the penalty (setting it to 1.0) moved AR-1.5B
Field-F1 from \arOneFiveFOneBefore\ to \arOneFiveFOneAfter.

Two properties of this defect deserve emphasis. First, it is
\textbf{model-size-dependent}: the penalty is \repPenaltySmall\ at 1.5B against
\repPenaltyLarge\ at 3B/7B, so it handicapped the smallest model most --- and
the 1.5B model is precisely the size-matched comparator in the claim under test.
Second, and more importantly, it was \textbf{one-sided}. The diffusion side was
never affected: \modelLLaDA\ carries no \texttt{repetition\_penalty} in its
configuration, and our denoising loop is hand-written and applies no logits
processors at all. There was no corresponding default waiting to handicap the
diffusion model, and so no possibility that the two errors would cancel. A
defect that can only degrade one arm of a two-arm comparison is not noise; it is
a bias with a known sign.

The defect surfaced only because two AR stacks disagreed. vLLM's
\texttt{SamplingParams} defaults \texttt{repetition\_penalty} to 1.0, so
migrating the AR baseline to a served runtime produced a Field-F1 gap on
identical models, prompts, events and (nominally) identical greedy decoding.
That gap has no legitimate explanation, and chasing it is what found the
penalty; with the penalty neutralised the two stacks agree to within ordinary
BF16 kernel drift. The generalizable check is therefore cheap: \emph{run the
baseline on two independent stacks and require the quality numbers to agree
before you compare it to anything}. We now pass \texttt{repetition\_penalty} as
an explicit constructor argument defaulting to 1.0, so that the decoding
configuration is an experimental choice on the record rather than whatever a
given checkpoint happens to ship.

\subsection{Defect 2: an unserved baseline}
\label{sec:defect2}

The second defect is coarser. Every AR latency we had measured came from a
naive HuggingFace \texttt{generate()} loop: one event at a time, no paged
attention, no continuous batching, no prefix reuse. The preprint under test
serves \emph{its} AR baselines with PagedAttention and FlashAttention-2, and
explicitly credits them with a frozen-conditioning cache, precisely so that its
diffusion model does not win on serving infrastructure. Our AR side was
therefore not the baseline the claim was made against.

Re-running the identical models, weights, prompts and events under vLLM moves
AR-1.5B p50 from \arOneFivePFiftyHF~ms to \arOneFivePFiftyVLLM~ms, and AR-7B
from \arSevenPFiftyHF~ms to \arSevenPFifty~ms. Nothing about the architecture
changed. Our prompt is a long fixed preamble followed by the single log line
that varies --- structurally the same object as the preprint's frozen
conditioning block --- and vLLM's prefix caching is what realises it; the tail
percentiles move even further than the median, because the naive loop's tail is
dominated by scheduling artifacts that a served runtime does not have.

\subsection{The corrected comparison}

\begin{table}[t]
\centering
\small
\setlength{\tabcolsep}{5pt}
\caption{Streaming latency and quality at concurrency $m=1$, $n=2000$ BGL
events, 12-field schema, greedy decoding, H100 PCIe 80GB. The top block is the
corrected comparison. The second block is the same AR models under the naive
HuggingFace loop with the checkpoint's repetition penalty active
(\S\ref{sec:defect1}--\S\ref{sec:defect2}); it is reported only to show the size
of the correction and is superseded by the top block. The bottom row is
\emph{claimed} by the preprint under test; its weights are unreleased, we did
not run it, and it must not be read as a reproduction. The diffusion row runs on
HuggingFace transformers because no diffusion runtime was installable on this
host at the time (\S\ref{sec:mdm-caveat}); a served diffusion measurement on a
diffusion-native runtime appears later in the paper. Tail ratio is
p99.9/p50.}
\label{tab:streaming}
\begin{tabular}{llrrrrr}
\toprule
Serving & Model & p50 (ms) & p99 (ms) & p99.9 (ms) & p99.9/p50 & Field-F1 \\
\midrule
\multicolumn{7}{l}{\emph{Corrected AR baseline (vLLM 0.19.1)}}\\
vLLM & AR-1.5B & \arOneFivePFiftyVLLM & \arOneFivePNineNine & \arOneFivePNineNineNine & \arTailRatioLo$\times$ & \textbf{\arOneFiveFOneAfter} \\
vLLM & AR-3B   & \arThreePFifty & \needsnum{AR-3B p99} & \needsnum{AR-3B p99.9} & \arTailRatioHi$\times$ & \textbf{\arThreeFOne} \\
vLLM & AR-7B   & \arSevenPFifty & \needsnum{AR-7B p99} & \needsnum{AR-7B p99.9} & \needsnum{AR-7B tail} & \textbf{\arSevenFOne} \\
\midrule
\multicolumn{7}{l}{\emph{Diffusion reference}}\\
HF & LLaDA-8B, $K{=}8$ & \mdmPFifty & \needsnum{LLaDA p99} & \mdmPNineNineNine & \textbf{\mdmTailRatio$\times$} & \mdmFOne \\
\midrule
\multicolumn{7}{l}{\emph{Superseded --- naive HF loop, repetition penalty active}}\\
HF & AR-1.5B (v1) & \arOneFivePFiftyHF & \needsnum{v1 1.5B p99} & \needsnum{v1 1.5B p99.9} & \needsnum{v1 1.5B tail} & \arOneFiveFOneBefore \\
HF & AR-7B (v1)   & \arSevenPFiftyHF & \needsnum{v1 7B p99} & \needsnum{v1 7B p99.9} & \needsnum{v1 7B tail} & \needsnum{v1 7B F1} \\
\midrule
\multicolumn{7}{l}{\emph{Claimed by the preprint --- not reproduced, weights unreleased}}\\
\emph{claim} & \emph{FlowFill-MDM, $K{=}8$} & \needsnum{claim p50} & \needsnum{claim p99} & \needsnum{claim p99.9} & \needsnum{claim tail} & \needsnum{claim F1} \\
\bottomrule
\end{tabular}
\end{table}

Table~\ref{tab:streaming} gives the corrected picture. Three of the conclusions
we had drawn from the uncorrected baseline reverse.

\paragraph{Reversal 1: parity is refuted.}
The near-parity finding paired AR-1.5B at \arOneFiveFOneBefore\ with LLaDA-8B at
\mdmFOne. The AR half of that pair was the penalty-suppressed number. Corrected,
AR-1.5B reaches \arOneFiveFOneAfter\ against the diffusion model's \mdmFOne, a
gap of \parityGap\ --- an order of magnitude wider than the \paperParityGap\ the
preprint claims. The parity was a property of the handicap, not of the
architectures.

\paragraph{Reversal 2: AR quality is flat in scale.}
Corrected Field-F1 across 1.5B/3B/7B is \arOneFiveFOneAfter\ /
\arThreeFOne\ / \arSevenFOne: non-monotonic, spread \arScaleSpread. Schema-bound
extraction from a single log line does not reward AR scale, which is a
substantive finding in its own right --- it means a 1.5B model is the correct
comparator and the study does not need a 7B one. It also explains the earlier
result cleanly. Our uncorrected numbers showed 7B comfortably ahead of 1.5B, and
we had read that as scale helping; what it actually measured is that the
repetition penalty is stronger at 1.5B (\repPenaltySmall) than at 7B
(\repPenaltyLarge). The apparent scaling trend was a plot of the defect.

\paragraph{Reversal 3: the direction of the speed claim.}
Properly served AR-1.5B runs at \arOneFivePFiftyVLLM~ms p50 against the
diffusion model's \mdmPFifty~ms: AR is \arSpeedupOverMDM$\times$ faster. The
claim under test is \paperClaimedSpeedup$\times$ in the opposite direction. We
are not asserting that the preprint's own measurements are wrong --- we cannot
run its model --- but the comparative direction does not reproduce with a served
AR baseline, and our earlier agreement with it was an artifact of comparing a
production runtime's absence against a hand-written denoising loop.

\subsection{What survives the correction}

The correction is not symmetric across claim types, and the distinction matters
more than any individual number. \textbf{Diffusion-internal} claims --- claims
about how the diffusion model behaves as a function of its own parameters ---
are untouched by anything in this section, because no AR measurement enters
them. \textbf{Comparative} claims are exactly the ones that fell.

The clearest survivor is the tail. LLaDA-8B has the tightest tail in
Table~\ref{tab:streaming} by a wide margin: p99.9/p50 of
\mdmTailRatio$\times$, against \arTailRatioLo--\arTailRatioHi$\times$ for the
served AR models. This follows from the architecture and not from the
comparison: a fixed $K$-step denoising schedule does the same work on every
input, so latency is content-independent by construction, while an AR model's
latency is a function of how many tokens the particular input makes it emit.
Content-independence is a real property of the diffusion model, it was measured
before the correction, and it is still there afterwards. What did not survive is
the \emph{bundle} --- the absolute speed claim that had been packaged together
with the bounded-tail claim. The tail ratio and the median are different
quantities, and only one of them was ever a diffusion property.

\subsection{The caveat on the other side}
\label{sec:mdm-caveat}

Symmetry requires stating the obvious objection: having just argued that an
unserved AR baseline invalidates a comparison, we present a diffusion row that
is itself unserved. That is a real limitation of Table~\ref{tab:streaming} and
we do not paper over it --- on this host, at that time, no production runtime
supported a masked diffusion model at all. The rest of the paper is in part a
response to it: the cost laws are measured with both architectures on the same
served stack, and the diffusion side is additionally measured on a
diffusion-native runtime, where its best quality (\dinferFOne\ Field-F1) comes
at \dinferPFifty~ms p50 --- slower still, in a cache-less configuration whose
limitations we document rather than a faster one. No configuration we can run
puts the diffusion model ahead of a served AR baseline on median latency for
this task at this output length.

\subsection{The lesson}

Both of our defects ran in the same direction. Neither was a mistake of
reasoning; both were defaults --- a value in a checkpoint's configuration file
that nobody chose, and a serving stack that was never selected because it was
never considered part of the experiment. Neither was visible in any result
table, and each one alone was sufficient to reverse the paper's headline
comparison.

The prescription follows directly. When the claim under test is ``architecture
X beats autoregression,'' the AR baseline is not the control, it is the
measurement --- price it first, and price it with the same care as the system
being promoted. Concretely: read the generation configuration of every
checkpoint you load and set every decoding parameter explicitly; serve the
baseline on the same class of runtime as the system it is being compared
against, and report the runtime; run the baseline on two independent stacks and
reconcile any quality difference before drawing a comparison; and check whether
each defect you find could have affected both arms or only one, because
one-sided defects do not average out.

We cannot say how often this happens in published comparisons; we have no
visibility into anyone's baseline but our own. What we can say is that we looked
carefully at ours, found two defects, and both of them favoured the new
architecture.
